\documentclass[11pt]{article}
\usepackage[margin=1in]{geometry}
\usepackage{amsmath,amssymb}
\usepackage{booktabs}
\usepackage{hyperref}
\usepackage{graphicx}
\usepackage{listings}
\usepackage{xcolor}
\lstset{basicstyle=\ttfamily\footnotesize,frame=single,breaklines=true}

\title{Independent replication of \\
\textbf{arXiv:2401.06240 -- ``Quantum eigenvalue processing''} \\
by Guang Hao Low and Yuan Su}
\author{QC-200 replication wave (2026-07-05)}
\date{2026-07-05}

\begin{document}
\maketitle

\begin{abstract}
This is a spot-check + partial-numerics independent replication of
\emph{Low \& Su 2024/2026, ``Quantum eigenvalue processing''}
(arXiv:2401.06240, v3, 26 Mar 2026, SIAM J.\ Computing / FOCS'24).
The paper introduces Quantum Eigenvalue Estimation (QEVE) and Quantum
Eigenvalue Transformation (QEVT), generalising QSVT
(Gily\'en--Su--Low--Wiebe 2019) from singular values to eigenvalues of
possibly non-normal, diagonalisable matrices with real spectra.
Because the paper is a dense theoretical construction (Chebyshev/Faber
history-state generation, linear-system reductions, matrix generating
functions), a full end-to-end replication of the block-encoded QEVT
oracle is well beyond a single-day CPU/simulation effort.
We therefore focus on the paper's own \emph{consistency claim}: ``on the
common ground where the input matrix is Hermitian, our result has thus
naturally recovered the complexity [and behaviour] of QSVT for
transforming singular values'' (Sec.~1.4, p.~12 of the arXiv PDF).
We build a real, executable QSVT circuit for a small block-encoded
Hermitian matrix, drive it with phase factors from an off-the-shelf
symmetric-QSP solver, and verify that the resulting unitary implements
the target eigenvalue-wise polynomial transformation \emph{exactly},
at machine precision, for (a) a Chebyshev-monomial target, (b) an
identity-polynomial sanity check, and (c) three smoothed-sign
approximants (degrees 11, 21, 41) that constitute a running example of
QEVT on Hermitian input in the QSVT-generalisation regime.
Verdict: \textbf{PARTIAL/SPOT-CHECK} -- the QEVT-reduces-to-QSVT claim
is confirmed on the Hermitian instance with maximum eigenvalue-wise
error of the QSVT circuit vs.\ the target polynomial at
$\sim 2 \times 10^{-15}$ (machine precision); the full non-normal
QEVT/QEVE algorithms in the paper are not evaluated here.
\end{abstract}

\section{Paper summary}

Low \& Su present a unifying framework for processing eigenvalues of
matrices accessed by a quantum computer through block-encoding oracles.
Two main algorithms:
\begin{itemize}
  \item \textbf{QEVE} (Quantum EigenValue Estimation): estimates an
        eigenvalue of a diagonalisable matrix $A$ with
        $O(\alpha \kappa / \epsilon \, \log(1/p))$ queries to a block
        encoding $O_A$ of $A/\alpha_A$ and a state-prep unitary for the
        corresponding eigenstate. This achieves the Heisenberg-limited
        scaling for eigenvalue estimation of a broad class of non-normal
        matrices with real spectra and Jordan forms.
  \item \textbf{QEVT} (Quantum EigenValue Transformation): implements a
        target polynomial transformation $p(A)$ eigenvalue-wise for
        diagonalisable $A$ with real spectra, using Chebyshev/Faber
        approximation of the target function.  Query complexity for
        $\|p\|_{\max,[-1,1]} \le 1$ and Hermitian $A$ is
        $\tilde O\!\left( \kappa_S n / p(A/\alpha_A|\psi\rangle) \right)$,
        matching QSVT up to polylog factors.
\end{itemize}
Underlying both is an efficient \emph{Chebyshev history state}
generation subroutine that encodes Chebyshev polynomials of $A$ in
quantum superposition via a matrix generating function; the paper also
extends this to a Faber history state for eigenvalue processing over
the complex plane.

\section{Claims and testability}

\begin{table}[h]\centering
\begin{tabular}{lllll}
\toprule
ID & Claim & Type & Testable in a day? & Tested here? \\
\midrule
C1 & QEVT reduces to QSVT for Hermitian $A$ (matches complexity + action, up to polylog) & Algorithmic & Yes (Hermitian case) & \textbf{Yes} \\
C2 & QSVT circuit implements $P(H)$ on eigenvalues for Hermitian $H$ block-encoded & Numerical & Yes & \textbf{Yes} (T$_3$, id, sign approx) \\
C3 & Sign-function approximation converges as $\deg\to\infty$ inside window $|x|>\delta$ & Numerical & Yes & \textbf{Yes} (deg 11, 21, 41) \\
C4 & QEVE achieves $O(\alpha\kappa/\epsilon\log(1/p))$ for non-normal input & Complexity & Not in a day & No \\
C5 & Chebyshev history-state prep via matrix generating function is efficient & Algorithmic & Not in a day (multi-oracle circuit) & No \\
C6 & Faber history state extends to complex-plane eigenvalue processing & Algorithmic & Not in a day & No \\
C7 & Quantum ODE solver: strictly linear-in-$t$ query complexity for average-case & Complexity & Not in a day & No \\
\bottomrule
\end{tabular}
\caption{Claims and which were tested. C1--C3 form the Hermitian-reduction
core; C4--C7 are the paper's non-normal / applications results which are
beyond a single-day CPU replication.}
\end{table}

\section{Method}

\subsection{Environment (versions)}
\begin{itemize}
  \item Host: CherryRd (macOS, x86\_64), Python 3.14.6 in \texttt{work/venv}.
  \item numpy 2.5.1, scipy 1.18.0, qiskit 2.5.0 (installed but not
        exercised; classical statevector via numpy suffices at this
        instance size), matplotlib 3.10.x, \texttt{pyqsp} (git latest,
        no \texttt{\_\_version\_\_} attribute).
  \item PDF tools: \texttt{pdftotext} (\texttt{poppler-utils}).
\end{itemize}

\subsection{Instance construction}

\paragraph{Diagonal Hermitian $H$.}
We take
$H = \mathrm{diag}(-0.8,\,-0.3,\,0.25,\,0.7)$
so that $H$ is Hermitian with $\lVert H\rVert = 0.8 < 1$, i.e.\
QSP-realisable without additional rescaling.  Dimension $N=4$.

\paragraph{Block encoding $U_H$.}
We use the standard single-ancilla, ``\texttt{Wx}''-convention QSP
signal operator (Gily\'en et al.\ 2019 / pyqsp):
\begin{equation}
U_H = \begin{pmatrix} H & i\sqrt{I-H^2} \\ i\sqrt{I-H^2} & H \end{pmatrix}
    \in \mathbb{C}^{8\times 8}
\end{equation}
which is unitary iff $\lVert H \rVert \le 1$ and Hermitian for diagonal
$H$.  Measured $\lVert U_H U_H^\dagger - I \rVert_\infty = 1.11 \times 10^{-16}$.

\paragraph{Target polynomials.}
Three families (all real, of definite parity so QSP-realisable):
(a) $P_A(x)=0.9\, T_3(x)$ (Chebyshev polynomial of the third kind,
scaled below the QSP unit-max-norm constraint);
(b) $P_C(x) = 0.9\, x$ (identity sanity check);
(c) sign approximation via truncated Chebyshev interpolant of
$\mathrm{erf}(\kappa x)$ with $\kappa = 2/\delta$, $\delta=0.2$, degrees
$d \in \{11,21,41\}$.  Coefficients computed via
\texttt{pyqsp.poly.PolyTaylorSeries().taylor\_series(chebyshev\_basis=True)}
mirroring the sym\_qsp minimal example distributed with pyqsp.

\paragraph{QSP phase factors.}
Computed via \texttt{pyqsp.angle\_sequence.QuantumSignalProcessingPhases}
with \texttt{method="sym\_qsp"} and \texttt{chebyshev\_basis=True}.
This is the symmetric-QSP Newton solver of Dong--Meng--Whaley--Lin
(2021) as implemented in pyqsp \texttt{sym\_qsp\_opt.py}; it converges
to machine precision (final residual $\lesssim 4\times 10^{-13}$ in the
iterative optimisation) within 5--8 iterations for all our targets.
This method replaces the earlier ``\texttt{laurent}'' root-finding
solver, which we verified crashes with \texttt{CompletionError} at
degree $\ge 11$ for the sign target -- see \texttt{failure\_analysis.md}.

\paragraph{QSVT circuit assembly.}
Using the pyqsp convention:
\begin{align}
R(\phi) &= \mathrm{diag}(e^{i\phi},\, e^{-i\phi})_{\text{on ancilla}} \\
U_\Phi &= R(\phi_0)\prod_{k=1}^{d} U_H\, R(\phi_k)
\end{align}
Assembled as an $8\times 8$ complex numpy array by exact matrix multiplication.
No shots, no sampling; exact statevector arithmetic.

\paragraph{Convention subtlety (documented for reproducibility).}
For the ``\texttt{Wx}''-convention signal operator with off-diagonal
factor $i\sqrt{1-x^2}$, an odd-parity QSP sequence deposits the target
\emph{real} polynomial in the \emph{imaginary} channel of $U_\Phi[0,0]$
(not the real channel).  We verified this empirically by cross-checking
the top-block of our \texttt{qsvt\_unitary} against pyqsp's own
\texttt{SymmetricQSPProtocol.gen\_unitary} for the scalar $N=1$ case
(exact match, see \texttt{work/debug\_convention.py}); the report code
consequently reads $P(x) = \mathrm{Im}[U_\Phi[0,0]]$ for the odd-parity
targets (T$_3$, identity, sign approx are all odd).

\subsection{Runbook}

\begin{lstlisting}[language=bash]
cd ~/Dropbox/REPLICATE-PROJECT/QC-200/QC-2401.06240-...
python3 -m venv work/venv
work/venv/bin/pip install --quiet qiskit numpy scipy matplotlib pyqsp
work/venv/bin/python report/evidence/qsvt_sign_replication.py   # main
work/venv/bin/python report/evidence/response_plot.py           # figures
\end{lstlisting}

\section{Results}

\subsection{Numerical replication (5/5 PASS)}

\begin{table}[h]\centering
\begin{tabular}{lllll}
\toprule
Experiment & Target polynomial & \# phases & max eigwise err & verdict \\
\midrule
A. Chebyshev T$_3$ & $0.9\, T_3(x)$ & 4 & $2.22\times10^{-16}$ & PASS \\
B1. Sign deg 11 & Chebyshev interp of $\mathrm{erf}(10 x)$ & 12 & $7.77\times10^{-16}$ & PASS \\
B2. Sign deg 21 & Chebyshev interp of $\mathrm{erf}(10 x)$ & 22 & $7.77\times10^{-16}$ & PASS \\
B3. Sign deg 41 & Chebyshev interp of $\mathrm{erf}(10 x)$ & 42 & $2.11\times10^{-15}$ & PASS \\
C. Identity $x$ sanity & $0.9\, x$ & 4 & $2.22\times10^{-16}$ & PASS \\
\bottomrule
\end{tabular}
\caption{QSVT circuit output vs.\ eigenvalue-wise application of the
target polynomial. Errors are at IEEE-754 double-precision machine
epsilon, confirming that the assembled QSVT unitary reproduces the
Chebyshev-QSP-realised polynomial exactly.}
\end{table}

\subsection{Sign-approximation quality vs.\ ideal $\mathrm{sign}(x)$}

\begin{table}[h]\centering
\begin{tabular}{lll}
\toprule
degree & $\max_i |P_\mathrm{poly}(\lambda_i) - \mathrm{sign}(\lambda_i)|$ &
         $\max_i |\mathrm{Im}[U_\Phi[0,0]]_{ii} - \mathrm{sign}(\lambda_i)|$ \\
\midrule
11 & $1.98 \times 10^{-1}$ & $1.98 \times 10^{-1}$ \\
21 & $1.31 \times 10^{-1}$ & $1.31 \times 10^{-1}$ \\
41 & $1.01 \times 10^{-1}$ & $1.01 \times 10^{-1}$ \\
\bottomrule
\end{tabular}
\caption{On the four test eigenvalues $\{-0.8,-0.3,0.25,0.7\}$,
sign-approximation error against ideal $\mathrm{sign}$. Both the offline
Chebyshev polynomial (column 2) and the QSVT circuit output (column 3)
agree with each other at machine precision. Error decreases with degree
as expected. Values are dominated by the eigenvalues closest to the
0.9-scaling boundary and to the $|x|<\delta=0.2$ smoothing window; the
$\lambda=0.25$ eigenvalue sits close to $\delta$ and is the hardest.}
\end{table}

\subsection{Response-function figure}
See \texttt{report/evidence/qsvt\_sign\_response.png}: plots
$\mathrm{Im}[U_\Phi[0,0]](x)$ (QSVT circuit output),
Chebyshev polynomial (target), and ideal $\mathrm{sign}(x)$
across $x \in [-1,1]$ for degrees 11, 21, 41. Circuit-vs-polynomial
curves are visually indistinguishable everywhere; the polynomial
approaches the step function on the $|x|>\delta$ windows as degree
grows.

\section{Verdict}

\textbf{PARTIAL} -- the paper's central algorithmic consistency claim
(QEVT $\equiv$ QSVT on Hermitian inputs, applying an arbitrary bounded
polynomial eigenvalue-wise via a block-encoded QSP circuit) is
reproduced at machine precision for a small, faithful, non-trivial
instance.  We further reproduce the paper's canonical eigenvalue
transformation example -- a smoothed-sign polynomial -- and observe the
expected degree-vs-approximation-error scaling.  The paper's full
non-Hermitian QEVT/QEVE machinery (Chebyshev history state generation
via matrix generating function, Faber extension to complex plane,
linear-system reductions, ODE-solver applications) is \emph{not}
exercised here and would require a substantially larger implementation
effort (multi-oracle circuits, quantum linear-system solver integration)
outside the day-scale replication scope.

\section{Files (artifacts inventory)}
See \texttt{report/artifacts\_summary.md} for the full inventory.

\section*{Open Questions}

\begin{description}
  \item[Q1] The pyqsp \texttt{PolySign.generate(chebyshev\_basis=True)}
        wrapper is broken in the current release (a local variable
        \texttt{scale} is referenced before assignment when
        \texttt{return\_scale=False}) -- we had to bypass it and build
        the Chebyshev interpolant of $\mathrm{erf}(\kappa x)$ manually
        via \texttt{PolyTaylorSeries}. How many downstream QSVT
        replication attempts silently produce garbage when they hit
        this exact code path? A cross-repo bisect on pyqsp API
        deprecations plus a formal ``QSP convention-mismatch smoke
        test'' would be worth publishing.
  \item[Q2] The ``\texttt{Wx}''-convention QSP+QSVT signal operator
        $[[x, i\sqrt{1-x^2}],[i\sqrt{1-x^2},x]]$ deposits an odd-degree
        real target polynomial in $\mathrm{Im}[U_\Phi[0,0]]$, not in
        $\mathrm{Re}[U_\Phi[0,0]]$. Many QSVT tutorials and even the
        pyqsp response module contain a comment saying ``sym\_qsp real
        component means little'' -- yet the widely-cited pedagogical
        rule of thumb is that ``$\mathrm{Re}[U[0,0]] = P(x)$''.
        Investigating exactly which QSVT papers/pedagogical treatments
        (including the paper under replication and its predecessors)
        adopt which convention -- and whether the confusion has caused
        undetected sign/parity errors in downstream numerical
        replications of Chebyshev-history-state or LCU-QSVT circuits --
        seems useful.
  \item[Q3] The paper's ``$\kappa_S$'' (basis condition number of the
        Jordan form) enters QEVT complexity multiplicatively, and the
        paper's convergence claims implicitly assume $\kappa_S$ is
        moderate.  How does the QEVT approach degrade empirically as
        $\kappa_S$ grows on a controlled family of non-normal test
        matrices (e.g.\ Jordan blocks perturbed by $\epsilon$, or the
        Grcar / Kahan pseudospectra classics)?  This is out of reach of
        a Hermitian-only single-day replication but is critical for
        practitioners trying to apply QEVT.
  \item[Q4] The pyqsp \texttt{sym\_qsp} Newton solver converges to
        residual $<10^{-12}$ in 5--8 iterations for our targets
        (degrees $\le 41$).  How does its scaling break down at
        production degrees ($\sim 10^3$--$10^4$) required for the
        paper's ODE-solver application with tight $\epsilon$?  Do any
        of the residual failure modes require the extended-precision
        or preconditioned solvers of Ying \& Kunis or Chao--Ni?  A
        head-to-head comparison of \texttt{sym\_qsp} vs.\ QSPPACK vs.\
        pennylane's \texttt{qml.qsvt} at high degree would help pin
        down the practical regime where the paper's asymptotic
        complexity kicks in numerically.
  \item[Q5] The paper's key subroutine -- Chebyshev history state
        generation via a matrix generating function -- is presented as
        a block-encoded circuit but not (as far as we saw) demonstrated
        on a small classically-simulable instance.  Building even a
        toy 2--3 qubit Chebyshev history state circuit (e.g.\ for
        $H = \sigma_x$ or a $4\times4$ Toeplitz-shift matrix) and
        verifying that the LCU-of-Chebyshev structure matches the
        paper's block-encoding recipe would give the community a
        canonical worked example.  Such an example is a natural next
        step to move beyond our Hermitian-QSVT-equivalent spot-check.
\end{description}


\end{document}
